\begin{proof}

\begin{enumerate}[label={$\langle1\rangle$\arabic*.},leftmargin=*]
\item Put $h=P(x,yz)$. It is weighted homogeneous of degree
$2d$ for the positive weights $(2,1,1)$.
On a small simply connected coordinate neighbourhood
take a holomorphic logarithm of the unit and set
\[
 \Phi(x,y,z)=\bigl(xu^{1/(md)},
                  yu^{1/(2md)},zu^{1/(2md)}\bigr).
                                                        \tag{A.BP.2}\hypertarget{tag-0050}{}\label{cdp:eq:A.BP.2}
\]
Its derivative at the origin is invertible, and direct
homogeneity gives $h\Phi=u^{1/m}h$. Thus
$f=(h\Phi)^m$ exactly. It suffices to construct a
compact spherical cycle of nonzero square in $h=1$,
because weighted homotheties then put it inside an
arbitrarily small actual fibre of \eqref{cdp:eq:A.BP.1}.

\item The curve $B=\{P(x,v)=1\}\subset\C^2$ is smooth:
Euler's identity $xP_x+vP_v=dP$ precludes a common
zero of its two partial derivatives on that level.
It is also connected. The map
\[
 B\longrightarrow\PP^1\setminus\{P=0\},
                       \quad(x,v)\longmapsto[x:v]
\]
is an etale cover of degree $d$. Locally a projective
representative has nonzero $P$-value, and its $d$
rescalings with $P=1$ give the sheets. Around any
one of the simple projective zeros of $P$, that
value has a simple zero, so its inverse $d$th root
has monodromy a full $d$-cycle. This monodromy is
transitive on the sheets, proving connectedness.

\item The function $v$ on $B$ has exactly $d$ distinct
simple zeros, since
$P(x,0)=P(1,0)x^d$ and $P_x(x,0)=d/x$ at those
points. Choose an embedded arc $\gamma$ in $B$
joining two of its zeros and avoiding all other
zeros. On the interior choose a continuous argument
of $v(\gamma(t))$ and two functions $A(t),B_0(t)$
with $A B_0=v(\gamma(t))$ and
$|A|=|B_0|=|v(\gamma(t))|^{1/2}$.
The map
\[
 (t,\theta)\longmapsto
   \bigl(x(\gamma(t)),A(t)e^{\mathrm i\theta},
                        B_0(t)e^{-\mathrm i\theta}\bigr)
                                                        \tag{A.BP.3}\hypertarget{tag-0051}{}\label{cdp:eq:A.BP.3}
\]
lies in $h=1$ and collapses both boundary circles
of $[0,1]\times S^1$ to points. It is therefore a
spherical cycle, with embedded interior circles.

\item A parallel-arc isotopy with the same endpoints gives
a homologous spherical cycle whose support meets
the first only at those endpoints. At each endpoint,
$P_x\ne0$, so $(y,z)$ are actual holomorphic
coordinates on $h=1$. The links of the two cycles
are parallel curves of slope $(1,-1)$ on a Clifford
torus in the local three-sphere. Their linking
number is $-1$. As in the explicit relative-chain
comparison in \ref{cdp:st:A.5.15}, that number is the local
intersection of the two coned cycles. Both
orientations reverse at the second endpoint,
so its intersection sign is unchanged. The
spherical class consequently has square $-2$.

\item The weighted homothety
$(x,y,z)\mapsto(\lambda^2x,\lambda y,\lambda z)$
maps $h=1$ biholomorphically to $h=\lambda^{2d}$
and shrinks the compact cycles and their isotopy
to the origin. Together with \eqref{cdp:eq:A.BP.2}, this puts
them inside an actual smooth global fibre of $f$.
That fibre is a complex two-torus and has zero
$\pi_2$, so the spherical class has zero homology
image. The oriented codimension-zero embedding
preserves the intersection of the compact cycles,
giving both zero and $-2$ for its square. This
contradiction proves the statement.

\item We now prove the variable-coefficient assertion,
including the boundary comparison. Absorb a local
$m$th root of the unit into $h$; this preserves the
form \eqref{cdp:eq:A.BP.4} and its two initial-form conditions.
For a complex parameter $s$ set
\[
 h_s=\sum_{i=0}^d a_i(s^2x,sy,sz)x^{d-i}(yz)^i.
                                                        \tag{A.BP.5}\hypertarget{tag-0053}{}\label{cdp:eq:A.BP.5}
\]
This is holomorphic also at $s=0$, with
$h_0=P_0(x,yz)$. For $s\ne0$ there is the exact
identity
\[
       h(s^2x,sy,sz)=s^{2d}h_s(x,y,z).           \tag{A.BP.6}\hypertarget{tag-0054}{}\label{cdp:eq:A.BP.6}
\]
The leading coefficient $a_0$ is a unit. The implicit
function theorem applied to its degree-$d$ root
polynomial gives distinct holomorphic functions
$b_i(s,x,y,z)$ and the exact factorization
\[
       h_s=a_0(s^2x,sy,sz)
                     \prod_{i=1}^d(x-b_i(s,x,y,z)yz)
                                                        \tag{A.BP.7}\hypertarget{tag-0055}{}\label{cdp:eq:A.BP.7}
\]
on one small parameter and coordinate neighbourhood.
Each factor defines a smooth surface, since its
$x$ derivative is a unit. For $i\ne j$ the ideal
generated by the two factors equals $(x,yz)$:
their difference is the unit $b_j-b_i$ times $yz$.
Thus the zero level has precisely the same two
critical curves $D=(x,yz)=0$ for every $s$.

\item Choose a small closed ball $K$ in the coordinate
chart such that the two axes of $D$ meet its boundary
transversely and the zero-level surfaces away from
those axes are transverse to the boundary. At
$s=0$ this follows directly for the smooth graph
surfaces $x=b_i(0)yz$ and sufficiently small radii.
It persists for small $s$ by compactness away from
neighbourhoods of the two boundary links.
Near either link, use normal coordinates $(U,V)$
for the smooth critical curve and a parameter $w$
along it. The normal initial term of $h_s$ is a
degree-$d$ binary form with distinct roots.
On the compact parameterized link its normal
gradient therefore satisfies a uniform bound
\[
 \|\partial_{U,V}h_s\|\ge c\|(U,V)\|^{d-1},
 \qquad |\partial_w h_s|\le C\|(U,V)\|^d
                                                        \tag{A.BP.8}\hypertarget{tag-0056}{}\label{cdp:eq:A.BP.8}
\]
after the neighbourhood is shrunk. The first bound
follows because the gradient of a homogeneous
reduced binary form has no zero on the unit sphere;
higher normal terms have one additional order.
The second follows from $h_s\in(U,V)^d$.

\item Transversality of the critical curve to $\partial K$
allows a boundary coordinate with one real component
of $w$ solved as a smooth function of $(U,V)$ and
the remaining component. The correction to normal
differentiation tangent to this boundary is a
bounded multiple of $\partial_w h_s$.
By \eqref{cdp:eq:A.BP.8} it is smaller by one normal order than
the surjective complex normal derivative. Thus
$h_s|_{\partial K}$ is a real submersion at every
nearby point with $h_s\ne0$, uniformly in $s$.
Away from the two link neighbourhoods, compactness
and the preceding zero-level transversality give
the same conclusion for all sufficiently small
nonzero values.

\item Remove the proper critical-value discriminant from
the product of the parameter disc and a sufficiently
small punctured value disc. The critical locus over
that domain cannot meet the boundary, so its
projection is proper and its image is complex
analytic. Its complement is connected. The level
sets in $K$ form a proper submersion with boundary
there, and Ehresmann transport identifies their
compact fibres by oriented diffeomorphisms.
Choose a small sphere cycle of the already proved
model $h_0=P_0(x,yz)$, and transport it to a fibre
of $h_{s_0}$ for $s_0\ne0$. Equation \eqref{cdp:eq:A.BP.6}
then embeds that fibre into the actual fibre
of $h^m$ with value $s_0^{2dm}\epsilon^m$.
It carries the spherical class of square $-2$
into a global torus. The same intersection and
asphericity contradiction proves \eqref{cdp:eq:A.BP.4}.

\end{enumerate}
\end{proof}
